\PassOptionsToPackage{dvipsnames}{color}
\PassOptionsToPackage{dvipsnames}{xcolor}
\documentclass[10pt]{NSP1}
\usepackage{url,floatflt}
\usepackage{helvet,times}
\usepackage{psfig,graphics}
\usepackage{mathptmx,amsmath,amssymb,bm}
\usepackage{float}
\usepackage[bf,hypcap]{caption}
%% --- additional packages required by the tables and figures of this paper ---
\usepackage{graphicx}
\usepackage{longtable,booktabs,array,calc}
\usepackage{multirow}

\providecommand{\tightlist}{\setlength{\itemsep}{0pt}\setlength{\parskip}{0pt}}
\providecommand{\real}[1]{#1}
\providecommand{\ul}[1]{\underline{#1}}
\providecommand{\hl}[1]{#1}
\providecommand{\RL}[1]{#1}
\providecommand{\textquotesingle}{'}
\DeclareUnicodeCharacter{03B1}{$\alpha$}
\DeclareUnicodeCharacter{03B2}{$\beta$}
\DeclareUnicodeCharacter{03C7}{$\chi$}
\DeclareUnicodeCharacter{03C1}{$\rho$}
\DeclareUnicodeCharacter{03BC}{$\mu$}
\DeclareUnicodeCharacter{03C3}{$\sigma$}
\DeclareUnicodeCharacter{2264}{$\leq$}
\DeclareUnicodeCharacter{2265}{$\geq$}
\DeclareUnicodeCharacter{2260}{$\neq$}
\DeclareUnicodeCharacter{00D7}{$\times$}
\DeclareUnicodeCharacter{00B1}{$\pm$}
\DeclareUnicodeCharacter{2248}{$\approx$}
\DeclareUnicodeCharacter{221A}{$\sqrt{\ }$}
\DeclareUnicodeCharacter{2192}{$\rightarrow$}
\DeclareUnicodeCharacter{2022}{$\bullet$}
\DeclareUnicodeCharacter{00B7}{$\cdot$}
\DeclareUnicodeCharacter{2032}{$'$}
\DeclareUnicodeCharacter{0393}{$\Gamma$}
\DeclareUnicodeCharacter{03B7}{$\eta$}
\DeclareUnicodeCharacter{03B8}{$\theta$}
\DeclareUnicodeCharacter{03BB}{$\lambda$}
\DeclareUnicodeCharacter{03A3}{$\Sigma$}
\DeclareUnicodeCharacter{03B4}{$\delta$}
\DeclareUnicodeCharacter{03B3}{$\gamma$}
\DeclareUnicodeCharacter{2026}{\ldots}
\DeclareUnicodeCharacter{FEFF}{}

\tolerance=1
\emergencystretch=\maxdimen
\hyphenpenalty=10000
\hbadness=10000

\topmargin=0.00cm

\def\sm{\smallskip}
\def\no{\noindent}

\def\firstpage{1}
\setcounter{page}{\firstpage}
\def\thevol{--}
\def\thenumber{--}
\def\theyear{20--}

\begin{document}

\titlefigurecaption{{\large \bf \rm Journal of Statistics Applications \& Probability }\\ {\it\small An International Journal}}

\title{Assessing the Readiness for Blended Educational Supervision in Jordan: Educational Supervisors' Perspectives from Irbid Governorate}

\author{Suad Eyadat$^{1}$, Najwa Abdel Hamid Darawsheh$^{2}$, Hanadi Alrashdan$^{3}$, Ekhlas Mohammad Abdel-Ghani ALRABA'I$^{4}$, Ayah Sayyah Ibrahim Alshamali$^{5}$, Kamal Subhi Said Nazzal$^{6}$, Shaman Saleh Marie Maabreh$^{7}$ and Ra'edah Ali Marashdeh$^{8}$}
\institute{$^{1}$Department of Educational Foundations and Administration, Faculty of Educational Sciences, Jadara University, Irbid, Jordan, s.eiadat@jadara.edu.jo\\
$^{2}$Department of Educational Foundations and Administration, Faculty of Educational Sciences, Jadara University, Irbid, Jordan, najwadarawsha@gmail.com, ORCID 0000-0002-0907-2188\\
$^{3}$Department of Educational Foundations and Administration, Faculty of Educational Sciences, Jadara University, Irbid, Jordan, h.alrashdan@jadara.edu.jo, ORCID 0000-0001-8898-3117\\
$^{4}$Department of Educational Foundations and Administration, Faculty of Educational Sciences, Jadara University, Irbid, Jordan, ekhlas.r@jadara.edu.jo, ORCID 0009-0008-6831-1688\\
$^{5}$Department of Educational Administration, Jadara University, Irbid, Jordan, a.shamali@jadara.edu.jo, ORCID 0009-0004-0035-1799\\
$^{6}$Department of Psychological Counseling, Faculty of Educational Sciences, Jadara University, Irbid, Jordan, Kamalnz@jadara.edu.jo\\
$^{7}$Department of Arabic Language and Literature, Faculty of Arts and Languages, Jadara University, Irbid, Jordan, Shaman.m@jadara.edu.jo\\
$^{8}$Department of Arabic Language and Literature, Faculty of Arts and Languages, Jadara University, Irbid, Jordan, raedaaali@jadara.edu.jo, ORCID 0009-0000-6520-1389}

\titlerunning{Readiness for blended educational supervision in Jordan}
\authorrunning{S. Eyadat et al.}

%corresponding author email
\mail{s.eiadat@jadara.edu.jo}

\received{7 Jun. 2026}
\revised{20 Jun. 2026}
\accepted{25 Jul. 2026}
\published{1 Aug. 2026}

\abstracttext{This study investigates the presence of conditions necessary for the application of the concept of blended educational supervision taking into account the views of educational supervisors. It meets the demand for supervision methods which allow enterprises to use the opportunities of digital transformation efficiently as well as to enhance the professional development of teachers. The study employs a descriptive-analytical method of research and uses a standard questionnaire for collecting data from a randomly selected sample of respondents. The questionnaire measures the following types of conditions: human resources, technological resources, and organizational resources. The data has been processed using both descriptive and inferential statistics. The analysis of data reveals that overall conditions are moderately present which means that blended educational supervision will only be partially implemented in the education system. The analysis indicates that the necessity of material and technological resources is much higher than that of human conditions. Most of the data had no statistically significant differences based on gender, experience, or education level, however, human conditions differ significantly in favor of the female respondents and those with PhD degree holders. In conclusion, it can be said that the effective implementation of blended educational supervision lies in enhancing the digital capacity, facilitating continuous professional growth, and providing institutional support. The study proposes to enlarge training options for supervisors, enhance their access to technological resources, and improve the organizational environment for blended supervision.}

\keywords{Digital transformation in education, instructional leadership, supervisory practices, professional development, school improvement}

\maketitle

\raggedbottom
\setlength{\textfloatsep}{12pt plus 2pt minus 2pt}
\setlength{\intextsep}{12pt plus 2pt minus 2pt}
\setlength{\floatsep}{12pt plus 2pt minus 2pt}

\section{Introduction}

The twenty-first century has ushered in a period of fast-tracked technological growth, enormous information diversification, and a leap towards knowledge-based and technology-oriented working models. These developments have had their impact on education systems worldwide, and this has not only been achieved through the implementation of different instruments, but also through the changes in the understanding of teaching quality, sustainability, institutional accountability, and school improvement worldwide \cite{ref1}. Educational supervision is significantly influenced by those changes, as it is states at the crossroads of educational policy and classroom practice. Traditionally, supervision was often associated with the idea of monitoring, obeying the rules, and occasional evaluations. Yet, today's education is increasingly changing the need for supervision to act as a surrounding process which helps the teachers organize their mentoring and professional dialogues and receive evidence-based feedback \cite{ref2,ref3,ref4}.

This turn to a new teaching and learning realities presents a major task for supervision procedures. Schools need to adjust in terms of digital transformation and introduction of online learning environments as well as data-based decision making. Meanwhile, supervisory practices are still guided mainly by traditional methods such as going into classrooms, making written reports and meeting face to face. Although all these practices are still useful, they fail to meet the challenges of contemporary teaching and learning. Supervisors need to assist teachers in their work both in real and online classrooms, give them feedback in a prompt manner, and use technology for preserving communication and professional follow-up \cite{ref5}.

As a result of these changes, blended educational supervision has been developed as a model that utilizes the advantages of direct supervision along with the tools of digital communication and online channels \cite{ref6}. This does not mean replacing direct observation of the classroom in any way or restricting supervision practice to purely electronic procedure. The primary purpose of blended supervision is to use classroom visits and professional dialogue, video conferencing and digital observation forms, recorded teaching materials, and online follow-up for the development of a single supervisory cycle. Therefore, it can be stated that blended supervision can help diminish both distance and time barriers, facilitate the communication between supervisors and teachers, and support reflective practice \cite{ref7,ref8,ref9}.

Nonetheless, it is essential to differentiate blended supervision from other similar terms that are often used interchangeably. Electronic supervision pertains to the use of electronic means such as email, e-learning platforms, or digital forms to perform supervision tasks remotely. On the other hand, digital supervision is more general as it refers to the application of digital technologies, systems, and environments for monitoring, feedback, communication, and learning. Blended supervision stands out through the combination of both direct and digital forms of supervision. In this case, the significance is attributed to the right ratio between in-person interaction and technology. Visiting the school provides the presence of important components of the process, face-to-face contacts, while electronic methods provide more opportunities for communication, documentation, feedback, and assistance in the process of supervision.

In educational environments like Jordan, the use of blended educational supervision is not consistent. The Ministry of Education has implemented some digital changes and modernization measures in educational services; however, the effectiveness of blended supervision needs more than just the availability of technology and internet access. The blended supervision should include proper regulations, qualified supervisors, the presence of digital competencies, administrative support, and established professional culture that is open to innovations \cite{ref10,ref11,ref12}. Very often, the greatest difficulty lies in the readiness of the supervision system to use the technology, and not in the fact that technology exists \cite{ref13}.

Global and regional academic literature shows an increase in interest for ICT-assisted supervision although important gaps remain. Several studies provided evidence of the effectiveness of digital tools in providing feedback, communication and support. However, insufficient clarity can be found in the literature concerning the methods of introduction of blended supervision in certain local systems with diverse infrastructure, legislation, training and culture.

In Jordan, the existing literature provided some useful information about electronic supervision and professional development but the data are still insufficient to form governorate-wise picture of the requirements for blended supervision from the perspective of the supervisors. Such need is especially relevant in Irbid Governorate where supervisory work depends on numerous school types and levels of technological readiness in maintaining the compliance of supervision activities with national transformations in the field of ICT.

As a result, the current research investigates to what extent the prerequisites for applying blended educational supervision are present in Irbid Governorate according to the viewpoint of educational supervisors. Additionally, the study provides suggestions for possible measures of improvement which might better facilitate the application of this form of supervision. The focus of the research is on the various kinds of required resources (human, technical, institutional) that should serve as a basis for providing proof of the effectiveness of realization of educational reform.

\subsection{Theoretical Background and Conceptual Framework}

Blended educational supervision is based on the belief that supervision must be a collaboration, development, and adaptation to teachers' professional needs. The authors \cite{ref14} argue that blended educational supervision is mainly characterized by participatory. Unlike the models which mainly focus on evaluation or conformity, blended educational supervision promotes cooperative inquiry, reflective conversations, and problem-solving among supervisory team, teachers and school authorities. Therefore, it is closer to developmental supervision than inspection since it concentrates on the learning of teachers and the improvement of teaching and learning.

According to the literature, blended supervision also depends on the proper balance of direct and indirect communication. Direct communication is the communication between teachers and supervisors on the face-to-face basis, including classroom visits, classroom observations, discussions and feedback. It is crucial communication at this stage since it provides supervisors with the broader picture of what is going on in the class; who communicate with whom and develop professional trust. Indirect communication, in its turn, is performed with the use of electronic and technological devices such as e-mail, virtual meetings, digital communication, recordings, and other feedback options. Indirect communication helps supervisors to communicate effectively with the teachers, provide guidance and track progress in the implementation of the recommendations made by them \cite{ref15}.

Blended supervision derives its merit from the combination of these two practices rather than relying on one more than the other. Direct supervision relates to building relational depth and providing contextual comprehension, while digital supervision implies continuity, flexibility, and gaining access to resources. When these two practices are combined, the supervision is transformed from being sporadic to continuous. The combination of these two forms of supervision can facilitate professional dialogue, contribute to improvements in feedback quality and enable teachers to reflect on their work overtime \cite{ref16}. Nevertheless, it needs to be mentioned that research reviews indicate that the blended supervision is not efficient only due to the use of online tools. The effectiveness of blended supervision is determined by the extent to which these tools are integrated into specific practices and being backed up by the institution policy, as well as the supervisory skills of the supervisor.

Thus, various requirements are essential in achieving blended educational supervision. The legal requirements refer to the laws, regulations and official instructions that approve and regulate the use of blended supervision. Such policies are designed to specify the functions of supervisors, ethical norms in the use of digital communication, data privacy, and system of accountability \cite{ref12}. The administrative requirements are the organizational arrangements needed for implementation. It is important to have clear mechanisms of communication, support from the leaders, and time allocation for implementation, and coordination among different directorates and schools \cite{ref17}. The human requirements are the knowledge, attitudes and digital skills of teachers and supervisors in such areas as understanding the purpose of blended supervision, using online supervision platforms properly, ethical online interaction, and the importance of person-to-person interaction in the framework of blended supervision \cite{ref18}.

Technical and professional requirements refer to the structured upbringing of supervisors, principals and teachers for blended supervision. They include training in digital observation, remote feedback, virtual communication and using technology for educational purposes \cite{ref19}.The material and technical requirements refer to the resources which are needed for implementation. Among such resources one can mention hardware, software, internet connection, digital platforms and technical support.

Finally, it could be said that all the mentioned requirements suggest that blended supervision is no longer just another technique but a process of an institutional change which requires some kind of coherence between the policy, physical condition, professionals' competences, and the culture of the school. Thus, to understand how feasible it would be to implement such a strategy in particular educational settings (e.g., in Irbid Governorate) we explored if these enabling factors are met here and where they require development. This article focuses on analyzing the feasibility of implementing blended supervision in Jordan in a specific context as opposed to discussing potential general advantages of applying this model in education.

\subsection{Research Problem}

The educational supervision technique in Jordan enhances teaching and learning methods, assist teachers with their professional advancement, and help achieve national educational objectives. However, despite the stated goals, the practice of supervision meets lots of challenges that prevent it from fully addressing the ongoing developments in education. The challenges of the supervision practice are especially relevant considering the increasing trend toward the digitalization of schools and the need for supervision to help teachers in the process of evaluation, as well as to provide ongoing support in the form of guidance, feedback, and professional development.

According to the researcher's experience in the field of educational supervision, a lot of traditional forms of supervision are still being applied, such as classroom observations, written reports, routine circulars and monitoring. Even though these techniques are still being used in some cases, when limited by the absence of technology or constant dialogue, they may turn out to be not so effective.

Previous studies conducted in Arab countries and Jordan has proven the necessity of reforming the supervisory practice. Al-Hanai \cite{ref20} stated that traditional supervision in Arab education does not completely correspond to today's pedagogical and technological achievements. Some other studies state that blended supervision may help improve educating quality, communication, and professional development under the condition of having provided all necessary frameworks, training, and support\cite{ref21,ref22,ref23}. However, the extent of their availability in concrete Jordanian governorates is not studied enough.

Thus, this research focuses on the question of whether the supervisory system in Irbid Governorate is ready to implement blended educational supervision. In more detail, the research aims to find out whether the necessary human, technical, and institutional means of success of this form of supervision exist from the point of view of educational supervisors. Additionally, it strives to find out the means of development needed for the strong possibility of successful and constant implementation of the system.

Based on this problem, the study addresses the following research questions:

\begin{enumerate}
\def\labelenumi{\arabic{enumi}.}
\item
  To what extent are the requirements for implementing blended educational supervision available in Irbid Governorate from the perspective of educational supervisors?
\item
  Are there statistically significant differences at the level of α = 0.05 in supervisors' perceptions of the availability of blended supervision requirements according to gender, years of professional experience, or academic qualification?
\item
  What developmental measures can be proposed to enhance the implementation of blended educational supervision in Irbid Governorate from the perspective of educational supervisors?
\end{enumerate}

\subsection{Research Objectives}

This paper aims at measuring the Readiness of applying Blended Educational Supervision in Irbid governorate through identification of its Basic strengths as well as pointing out what needs more improvement. It also performs a comparative analysis between the perception of Male vs. Female Supervisors regarding the presence of Blended Supervision, making recommendation to be able to develop Blended educational supervision in Irbid Governorate.

\subsection{Significance of the Study}

This study is important because of its impact on the blended learning supervision in Jordan. By analyzing the presence of certain variables that determine the success of the introduction of blended supervision into the local educational practice, this study helps gain better insights into the concept of blended supervision as well as its implementation in educational systems of the area. At the same time, it offers a model which can help develop future research in the field of supervisory development, digital transformation, and professional learning in Arab states.

In practical terms, the study can be beneficial for educational policymakers, educational administrators, and education supervising departments. The findings can be used for determining the types of resources, capacities, and conditions necessary for successful blended supervision implementation. Additionally, the findings can help in determining the right digital solutions, training activities and organizational procedures supporting both face-to-face and online supervision.

\subsection{Conceptual and Operational Definitions}

Educational supervision refers to the process characterized by cooperation and the pursuit of advancement through which supervisors assist teachers and school managers with their guidance to accomplish better teaching results, \cite{ref24}. In this study, educational supervision will be understood as the types of supervisory acts conducted in the context of Irbid Governorate -- either directly or indirectly, for the purpose of facilitating improvement of teaching.

Blended educational supervision is a system of supervision that includes both direct face-to-face method and digital methods of supervision. This is achieved through the usage of information and communication technologies with an emphasis on the aspects of in-person contact between the supervisor and supervisor. In this study, the operational definition of blended educational supervision will refer to the use of both direct and indirect methods of supervision.


\section{Empirical Studies}
Recently, there has been a surge in empirical research into technology-based supervision in education to find more versatile solutions for aiding instructors and enhancing the quality of education. In different settings, the findings indicate that technologies based on electronic, digital, and hybrid forms are capable of improving communication, feedback, and professional reflection. Nevertheless, they state that implementation depends on the condition of the relevant infrastructures, the capabilities of teachers and supervisors with regard to technologies, and the existence of supporting policies.

The initial studies in the field focused mainly on the potential of using virtual supervision for communication and feedback enhancement. For instance, Schwartz-Bechet \cite{ref25} utilized qualitative research with a form of collecting data through semi-structured meetings, supervisors' notes, and reflective emails between supervisors and teachers in training. This research revealed that IT-based supervision supported formative feedback, prompted reflective thinking, and made it easier for the supervisors to support monitoring. Another example is the research conducted by Arnauld \cite{ref26}, who used a single case qualitative study with prospective teachers from two Southwestern states of the USA, focusing on their exchanges with their supervisors.

Research related to the availability of requirements for electronic or blended supervision in Arab educational contexts has been done. For instance, Hamdan \cite{ref27}, adopting a descriptive-analytical design, researched e-supervision readiness in a sample of 177 supervisors in public schools in Gaza. The results of the research showed moderate readiness and highlighted deficiencies in technological infrastructure and digital training. Hazaimeh \cite{ref18}arrived at similar conclusions from the descriptive survey of 233 supervisors in Northern Governorates of Jordan, which also rated the availability of e-supervision requirements as moderate. These studies agree that Arab supervisory systems have begun adopting technology-based methods, but their development is limited by weak infrastructure and lack of necessary professional training. However, these two studies differ in the context and scope. While Hamdan studied only public schools in Gaza, Hazaimeh was able to explore Northern Jordan without focusing on a particular governorate.

Moreover, some research has developed specific models for introducing change, such as Rajabi and Zamil's \cite{ref23} mixed-methods study involving 584 supervisors and school principals from Palestine. Quantitatively, they found that while states of readiness regarding administration were high, use of digital communication, organization infrastructure, supervisory skills, and involvement of stakeholders reached a moderate level only. In qualitative results, authors highlight the necessity to create trustworthy online spaces along with ensuring regular management follow-up. Another example of research implementation focus is the descriptive survey by Al-Juhani and Al-Juhani \cite{ref28} these researchers involved 278 supervisors and also engaged a considerable amount of teachers from Medina.

The research outcomes imply that compared with other countries, the organizations' readiness (organizational, human, and material) is higher in Saudi Arabia. Hence, these findings suggest that if there are appropriate institutional provisions and infrastructures, blended supervision could develop even more efficiently. As can be seen, differences between findings obtained in Palestine and Saudi Arabia prove how deep contextual knowledge of blended supervision exists within each country.

Regarding Professional Competence, according to Winn (2009), the efficiency of electronic and blended supervision critically relies on the readiness and digital competence of supervisors and teachers. In his later works, he added that using electronic platforms demands particular skills from participants that should be constantly improved via continuous professional development. Moreover, Banga and Al-Qathami \cite{ref29} argue that blended learning success highly relies on both teachers and supervisors' technologies skills and ethically applying digital tools. Salman \cite{ref22} adds that teachers and educators need to ensure continuous professional development in order to apply blended learning techniques effectively.

Recent research shows the significance of organizational culture and reflective collaboration. The researchers created a model of professional development supervision for novice teachers that included both communication and mediation. It was concluded that online supervision promotes professional growth if it is combined with reflective dialogue and decision-making. Similarly, Mohamed, Hassan, and Salim \cite{ref30} examined blended supervision in secondary schools of Azhar and found that organizational issues, lack of drive, and insufficient training were barriers to success.

Overall, the previous studies show that the technology-assisted supervision can be highly effective in enhancing feedback, communication, and reflective practice as long as they are implemented correctly. The studies exploring the situation in Arab countries share that the readiness of the organizations is rather moderate, especially as for the infrastructure, training, and human resources. The literature showed that there are important differences in methodology and context of the investigations. While some of the multidisciplinary studies are qualitative and small-scale, others are descriptive studies, and some of them represent mixed methods and models. However, need for more local evidence that assesses the readiness for blended supervision at the level of the governorate in Jordan still exists. This research addresses this gap due to the assessment of whether the human, technical, and institutional conditions for blended supervision are available in Irbid Governorate.

\subsection{Research Gap and Contribution}

While there have been studies on digital, electronic, and blended supervision, most studies focus on specific aspects of blended supervision, such as teacher readiness, attitudes, and the use of digital tools. A number of studies done in Gaza, Palestine, and Saudi Arabia provide useful data on how blended supervision can be effective and what factors support its implementation \cite{ref27,ref23,ref28}. However, these studies have been conducted in different educational systems that have different policies, resources, and institutions than those present in Jordan. Hence, findings obtained in these studies can't be directly generalized to the Jordanian context; it would require thorough consideration of the specific context of Jordan.

Existing studies in Jordan have generated some familiarity with e-supervision and supervision development, particularly in the Northern governorates \cite{ref18,ref22}. However, literature offers limited information on what the requirements of blended supervision mean at the level of specific governorates, which is important considering that the level of readiness for implementation can be affected by local infrastructure, training opportunities, administrative support, and supervisors' backgrounds. In this sense, Irbid Governorate is relevant when it comes to exploring how national goals of digital transformation are reflected on the ground.

As a result, the contribution of this research is primarily contextual and practical. The current research does not aim to institute an innovative method, since descriptive and inferential methods have already been employed in previous research. On the contrary, the significance lies in the governorate-level evidence about the necessity of human, technical, and institutional conditions for finding blended supervision in the context of Irbid. Furthermore, it ties the empirical findings to measures of development that can be valuable for the educational leaders and decision-makers in pointing out the gaps, especially those related to professional training, infrastructure development, and proper condition of institutions.

\section{Methodology and Procedures}

\subsection{Research Design}

In this study, the descriptive-analytical methodology has been employed. This methodology is best suited for investigation of educational phenomena that take place naturally \cite{ref31}. Our choice was determined by the purpose of the study, which was to investigate educational supervisors' views on having the necessary resources for implementing blended supervision in Irbid Governorate. This design allowed the author to additionally investigate the current state of readiness and development solutions based on participants' answers. The descriptive-analytical design is widely used for educational purposes wherein various perceptions, practices, or conditions are investigated \cite{ref32}.

\subsection{Population of the Study}

The study sample included all educational supervisors employed at the educational directorates present in Irbid Governorate during the academic year 2025. The total number of educational supervisors in the governorate was obtained from the official statistics released from the communications with the Jordanian Ministry of Education and the concerned educational directorates.

\subsection{Sample of the Study}

A simple random sampling technique was employed to select a sample of 180 educational supervisors from the target population. The study's sample size was considered adequate for the descriptive and inferential statistical analyses employed in the study, in accordance with established criteria for quantitative research in social sciences. Table~\ref{tab1} indicates the distribution of the participant sample according to their sex, experience, and education level.

{\small
\begin{longtable}[]{@{}
  >{\raggedright\arraybackslash}p{(\columnwidth - 6\tabcolsep) * \real{0.3367}} >{\raggedright\arraybackslash}p{(\columnwidth - 6\tabcolsep) * \real{0.1892}} >{\raggedright\arraybackslash}p{(\columnwidth - 6\tabcolsep) * \real{0.1972}} >{\raggedright\arraybackslash}p{(\columnwidth - 6\tabcolsep) * \real{0.2769}}@{}}
\caption{Frequencies and Percentages of the Study Sample According to Demographic Variables}
\label{tab1}\\
\toprule\noalign{}
\endhead
\bottomrule\noalign{}
\endlastfoot
\textbf{Variable} & \textbf{Category} & \textbf{Frequency} & \textbf{Percentage (\%)} \\
\textbf{Gender} & Male & 105 & 59.2 \\
& Female & 75 & 40.8 \\
& \textbf{Total} & \textbf{180} & \textbf{100.0} \\
\textbf{Experience} & Less than 10 years & 101 & 57.4 \\
& 10 years or more & 79 & 42.6 \\
& \textbf{Total} & \textbf{180} & \textbf{100.0} \\
\textbf{Educational Qualification} & Higher Diploma & 84 & 38.1 \\
& Master's & 52 & 33.6 \\
& Doctorate & 44 & 28.3 \\
& \textbf{Total} & \textbf{180} & \textbf{100.0} \\
\end{longtable}
}

\subsection{Instrument of the Study}

\noindent In order to reach the objectives of this study, a questionnaire was created to examine the existence of the requisites for practice of blended pedagogical supervision in the Irbid Governorate. The tool was constructed after a review of the literature about educational supervision and previous studies, including the works of \cite{ref29,ref22,ref28}, and \cite{ref30}. The questionnaire originally consisted of two parts. The first part measured the existence of the requisites for effecting the blended educational supervision, and consisted of 20 items divided into three domains: technical-electronic requisites, material requisites, and human requisites. The second part contained 19 items and focused on the developers' proposals for the study.

\noindent The developed draft questionnaire was evaluated by relevant experts in the field of educational supervision, educational measurement, and research methodology. Their recommendations revolved around the simplicity of items, the quality of wording, congruency with the study objectives, and the appropriateness of the study for the Jordanian educational setting. Modifications were made according to the comments of the experts who were at least 80\% in agreement with the provided feedback. Thus, after the evaluation, three items were omitted or combined together due to the overlap or lack of clarity of the items. The final version of the questionnaire contained 36 items in total.

\subsection{Validity of the Instrument}

The content validity of the scale was assessed through a group of experts. The group of experts reviewed the questionnaire items for relevance, clarity, and suitability with respect to the aims of the study and dimensions being measured. Their comments were utilized in order to refine the content of the questionnaire.

As a next step, the questionnaire was tested on 30 educational supervisors in order to determine the construct-related validity. Pearson correlation coefficients were then computed in order to analyze the correlations between dimensions of the instrument. The results are presented in Table 2.

{\small
\begin{longtable}[]{@{}
  >{\raggedright\arraybackslash}p{(\columnwidth - 12\tabcolsep) * \real{0.3141}} >{\raggedright\arraybackslash}p{(\columnwidth - 12\tabcolsep) * \real{0.1229}} >{\raggedright\arraybackslash}p{(\columnwidth - 12\tabcolsep) * \real{0.1074}} >{\raggedright\arraybackslash}p{(\columnwidth - 12\tabcolsep) * \real{0.1085}} >{\raggedright\arraybackslash}p{(\columnwidth - 12\tabcolsep) * \real{0.0950}} >{\raggedright\arraybackslash}p{(\columnwidth - 12\tabcolsep) * \real{0.1804}} >{\raggedright\arraybackslash}p{(\columnwidth - 12\tabcolsep) * \real{0.0716}}@{}}
\caption{Inter-correlation Coefficients among the Instrument Dimensions}
\label{tab2}\\
\toprule\noalign{}
\endhead
\bottomrule\noalign{}
\endlastfoot
\textbf{Relationship} & \textbf{Statistic} & \textbf{Technical-electronic} & \textbf{Material} & \textbf{Human} & \textbf{Developmental Proposals} & \textbf{Total Scale} \\
\textbf{Technical-electronic} & r & 1 & & & & \\
& Sig. & & & & & \\
\textbf{Material} & r & .655** & 1 & & & \\
& Sig. & .000 & & & & \\
\textbf{Human} & r & .346** & .345** & 1 & & \\
& Sig. & .000 & .000 & & & \\
\textbf{Developmental Proposals} & r & .211** & .237** & .630** & 1 & \\
& Sig. & .000 & .000 & .000 & & \\
\textbf{Total Scale} & r & .815** & .755** & .720** & .659** & 1 \\
& Sig. & .000 & .000 & .000 & .000 & \\
\end{longtable}
}

\textbf{*p \textless{} 0.05; *\emph{p \textless{} 0.01}}

\noindent As shown in Table~\ref{tab2}, the correlations among the sub-dimensions ranged from .211 to .655, while correlations between the dimensions and the total scale ranged from .659 to .815. These values indicate acceptable internal relationships among the dimensions and support the use of the questionnaire for the purposes of this study. However, no exploratory factor analysis or confirmatory factor analysis was conducted. Therefore, the evidence for the factorial structure of the instrument should be considered preliminary, and this issue is acknowledged as a methodological limitation of the study.

\subsection{Reliability of the Instrument}

To assess reliability, the test--retest method was used by giving the instrument a second time to a group of 25 supervisors (not part of the study sample) two weeks later. A high level of temporal stability was shown by the two measures' 0.92 Pearson correlation. Cronbach's alpha was used to test internal consistency, and the entire scale's reliability value of 0.86 was higher than the social sciences' generally accepted criterion of 0.70 (Taber, 2018). Table~\ref{tab3} displays the outcomes of both reliability tests.

{\small
\begin{longtable}[]{@{}
  >{\raggedright\arraybackslash}p{(\columnwidth - 4\tabcolsep) * \real{0.5080}} >{\raggedright\arraybackslash}p{(\columnwidth - 4\tabcolsep) * \real{0.2027}} >{\raggedright\arraybackslash}p{(\columnwidth - 4\tabcolsep) * \real{0.2893}}@{}}
\caption{Cronbach's Alpha and Test--retest Reliability Coefficients}
\label{tab3}\\
\toprule\noalign{}
\begin{minipage}[b]{\linewidth}\raggedright \textbf{Dimension} \end{minipage} & \begin{minipage}[b]{\linewidth}\raggedright \textbf{Test--retest} \end{minipage} & \begin{minipage}[b]{\linewidth}\raggedright \textbf{Cronbach's Alpha} \end{minipage} \\
\midrule\noalign{}
\endhead
\bottomrule\noalign{}
\endlastfoot
Technical-electronic requirements & 0.89 & 0.90 \\
Material requirements & 0.90 & 0.84 \\
Human requirements & 0.88 & 0.83 \\
Developmental proposals & 0.79 & 0.96 \\
\textbf{Overall Scale} & \textbf{0.86} & \textbf{0.88} \\
\end{longtable}
}

These coefficients indicate high levels of internal consistency and stability to confirm the reliability of the measurement instrument for research purposes.

\subsection{Data Collection Procedures}

After obtaining the necessary administrative approval from the relevant educational authorities, the questionnaire was distributed to the selected educational supervisors in Irbid Governorate. Participants were informed of the purpose of the study and were assured that their responses would be used only for academic research purposes. Participation was voluntary, and confidentiality was maintained throughout the data collection process. Completed questionnaires were assessed prior to data being housed in SPSS to ensure that the said data were fit for further analysis.

\subsection{Scoring and Interpretation Procedures}

The final tool consisted of 36 positively worded items broken down into four areas: technical-electronic requirements, physical requirements, human requirements, and measures of development. The responses were recorded using a 5-point Likert scale where 5 means ``strongly agree,'' 4 means ``agree,'' 3 means ``neutral'', 2 means ``disagree'' and 1 means ``strongly disagree.''

Interpretation of the mean scores entailed dividing the range of the scale into three equal parts: low, medium, and high in range. The width of the category was obtained by deducting the lowest score from the highest score and dividing the obtained result by the number of categories.

Participation was voluntary, and confidentiality was maintained throughout the data collection process. Completed questionnaires were reviewed to ensure that they were suitable for statistical analysis before the data were entered into SPSS.

To interpret the mean scores, the scale range was divided into three equal categories: low, moderate, and high. The category width was calculated by subtracting the minimum score from the maximum score and dividing the result by the number of categories, as follows:

Category width = (Maximum score -- Minimum score) / Number of categories\\
Category width = (5 -- 1) / 3 = 1.33

Accordingly, the interpretation criteria were as follows:

{\small
\begin{longtable}[]{@{}
  >{\raggedright\arraybackslash}p{(\columnwidth - 2\tabcolsep) * \real{0.5630}} >{\raggedright\arraybackslash}p{(\columnwidth - 2\tabcolsep) * \real{0.4370}}@{}}
\caption{Criteria for interpreting mean scores}
\label{tab4}\\
\toprule\noalign{}
\begin{minipage}[b]{\linewidth}\raggedright \textbf{Mean Range} \end{minipage} & \begin{minipage}[b]{\linewidth}\raggedright \textbf{Interpretation} \end{minipage} \\
\midrule\noalign{}
\endhead
\bottomrule\noalign{}
\endlastfoot
1.00 -- 2.33 & Low \\
2.34 -- 3.67 & Moderate \\
3.68 -- 5.00 & High \\
\end{longtable}
}

The category range (1.33) was calculated using the formula:

\begin{equation}
\frac{\text{Maximum} - \text{Minimum}}{\text{Number of Categories}} = \frac{5-1}{3} = 1.33 .
\end{equation}

\begin{equation}
\text{Interval} = \frac{\text{Max Score} - \text{Min Score}}{3} = \frac{5-1}{3} = 1.33 .
\end{equation}

This classification facilitated the interpretation of respondents' perceptions regarding the degree of availability of the blended supervision requirements and developmental proposals.

\subsection{Statistical Analysis}

The data gathered has been analyzed using SPSS software version 26. Descriptive statistical methods, which include the calculation of frequencies, percentages, and means, have been used to grasp the demographic characteristics of study participants and address research questions about the status of availability of blended supervision requirements and development proposals. The use of Pearson correlation coefficients allowed us to determine the relationships between the questionnaire dimensions and collect initial evidence of internal consistency among the domains. In addition, the internal reliability of the instrument was estimated through estimating Cronbach's alpha.

Inferential statistical methods were used to identify the differences in the participants' perception of supervision based on gender, professional experience, and level of education. A three-way ANOVA without interaction effects was conducted to analyze the differences in total score. Later, multivariate ANOVA was conducted in order to analyze the differences in the key dimensions of the questionnaire. The significance of the data was checked as α = .05.

The data were confirmed for completeness and readiness for analysis prior to running the inferential analyses. Random sampling of the data got ensured by the independence of observations. However, the use of the advanced testing of assumptions was not applied in the research, which is the limitation of the study.

\section{Results}

\subsection{Results Related to the First Research Question}

The first research question examined the extent to which the requirements for implementing blended educational supervision are available in Irbid Governorate from the perspective of educational supervisors. To answer this question, arithmetic means and standard deviations were calculated for each dimension of the questionnaire. Table~\ref{tab5} presents the results ranked in descending order according to the mean scores.

{\small
\begin{longtable}[]{@{}
  >{\raggedright\arraybackslash}p{(\columnwidth - 8\tabcolsep) * \real{0.2150}} >{\raggedright\arraybackslash}p{(\columnwidth - 8\tabcolsep) * \real{0.5191}} >{\raggedright\arraybackslash}p{(\columnwidth - 8\tabcolsep) * \real{0.0972}} >{\raggedright\arraybackslash}p{(\columnwidth - 8\tabcolsep) * \real{0.0726}} >{\raggedright\arraybackslash}p{(\columnwidth - 8\tabcolsep) * \real{0.0960}}@{}}
\caption{Arithmetic Means and Standard Deviations of the Degree of Availability of Blended Supervision Requirements in Irbid Governorate (N = 180)}
\label{tab5}\\
\toprule\noalign{}
\begin{minipage}[b]{\linewidth}\raggedright \textbf{Rank} \end{minipage} & \begin{minipage}[b]{\linewidth}\raggedright \textbf{Dimension} \end{minipage} & \begin{minipage}[b]{\linewidth}\raggedright \textbf{Mean} \end{minipage} & \begin{minipage}[b]{\linewidth}\raggedright \textbf{SD} \end{minipage} & \begin{minipage}[b]{\linewidth}\raggedright \textbf{Level} \end{minipage} \\
\midrule\noalign{}
\endhead
\bottomrule\noalign{}
\endlastfoot
1 & Material Requirements & 3.92 & 1.19 & High \\
2 & Technical--Electronic Requirements & 3.96 & 1.06 & High \\
3 & Human Requirements & 3.85 & 1.01 & High \\
\textbf{Overall Scale} & & \textbf{3.90} & \textbf{0.51} & \textbf{High} \\
\end{longtable}
}

As shown in Table~\ref{tab5}, the overall mean score for the availability of the requirements for implementing blended educational supervision was 3.90 (SD = 0.51), thus showing that the overall availability is at a high level.

Out of the three dimensions, Technical--Electronic Requirements ranked first, with a mean score \textbf{of} 3.96 \textbf{(SD =} 1.06\textbf{),} followed by Material Requirements, with a mean score of 3.92 (SD = 1.19). Both dimensions were classified as High. In contrast, human requirements ranked last, with a mean score of 3.85 (SD = 1.01), which falls within the High level according to the interpretation criteria adopted in this study.

\subsection{Results of Second Question}

The second research question examined whether there were statistically significant differences at α = .05 in educational supervisors' perceptions of the availability of blended educational supervision requirements according to gender, professional experience, and educational qualification. Means and standard deviations were first calculated for the total scale scores according to these variables. The results are presented in Table 6.

{\small
\begin{longtable}[]{@{}
  >{\raggedright\arraybackslash}p{(\columnwidth - 8\tabcolsep) * \real{0.3599}} >{\raggedright\arraybackslash}p{(\columnwidth - 8\tabcolsep) * \real{0.2541}} >{\raggedright\arraybackslash}p{(\columnwidth - 8\tabcolsep) * \real{0.1000}} >{\raggedright\arraybackslash}p{(\columnwidth - 8\tabcolsep) * \real{0.1748}} >{\raggedright\arraybackslash}p{(\columnwidth - 8\tabcolsep) * \real{0.1112}}@{}}
\caption{Means and Standard Deviations of Supervisors' Scores by Demographic Variables (N = 180)}
\label{tab6}\\
\toprule\noalign{}
\begin{minipage}[b]{\linewidth}\raggedright \textbf{Variable} \end{minipage} & \begin{minipage}[b]{\linewidth}\raggedright \textbf{Category} \end{minipage} & \begin{minipage}[b]{\linewidth}\raggedright \textbf{N} \end{minipage} & \begin{minipage}[b]{\linewidth}\raggedright \textbf{Mean} \end{minipage} & \begin{minipage}[b]{\linewidth}\raggedright \textbf{SD} \end{minipage} \\
\midrule\noalign{}
\endhead
\bottomrule\noalign{}
\endlastfoot
\textbf{Gender} & Male & 105 & 3.92 & 0.49 \\
& Female & 75 & 3.89 & 0.53 \\
\textbf{Experience} & Less than 10 years & 101 & 3.88 & 0.48 \\
& 10 years or more & 79 & 3.94 & 0.51 \\
\textbf{Educational Qualification} & Higher Diploma & 84 & 3.97 & 0.47 \\
& Master's & 52 & 3.88 & 0.50 \\
& Doctorate & 44 & 3.85 & 0.49 \\
\end{longtable}
}

As shown in Table~\ref{tab6}, small differences were obserbed in the mean scores according to gender, experience, and educational qualification. To test whether these differences were statistically significant, a three-way analysis of variance without interaction effects was conducted on the total scale score. The results are shown in Table 7.

{\small
\begin{longtable}[]{@{}
  >{\raggedright\arraybackslash}p{(\columnwidth - 10\tabcolsep) * \real{0.4056}} >{\raggedright\arraybackslash}p{(\columnwidth - 10\tabcolsep) * \real{0.1739}} >{\raggedright\arraybackslash}p{(\columnwidth - 10\tabcolsep) * \real{0.1039}} >{\raggedright\arraybackslash}p{(\columnwidth - 10\tabcolsep) * \real{0.1435}} >{\raggedright\arraybackslash}p{(\columnwidth - 10\tabcolsep) * \real{0.0852}} >{\raggedright\arraybackslash}p{(\columnwidth - 10\tabcolsep) * \real{0.0879}}@{}}
\caption{Three-Way ANOVA Results for Total Scale Scores}
\label{tab7}\\
\toprule\noalign{}
\begin{minipage}[b]{\linewidth}\raggedright \textbf{Source of Variance} \end{minipage} & \begin{minipage}[b]{\linewidth}\raggedright \textbf{SS} \end{minipage} & \begin{minipage}[b]{\linewidth}\raggedright \textbf{df} \end{minipage} & \begin{minipage}[b]{\linewidth}\raggedright \textbf{MS} \end{minipage} & \begin{minipage}[b]{\linewidth}\raggedright \textbf{F} \end{minipage} & \begin{minipage}[b]{\linewidth}\raggedright \textbf{Sig.} \end{minipage} \\
\midrule\noalign{}
\endhead
\bottomrule\noalign{}
\endlastfoot
Gender & 0.055 & 1 & 0.055 & 0.136 & 0.713 \\
Experience & 0.372 & 1 & 0.372 & 0.920 & 0.339 \\
Educational Level & 1.268 & 2 & 0.634 & 1.569 & 0.211 \\
Error & 70.736 & 175 & 0.404 & & \\
Total & 1742.693 & 179 &  &  &  \\
\end{longtable}
}

\begin{quote}
\emph{\textbf{p \textless{} 0.05}}
\end{quote}

As appeared in Table~\ref{tab7}, no statistically significant differences were evident in total scores based on gender or professional experience and Educational Level.

\subsection{Results of Third Question}

To answer the third research question \emph{``What are the developmental proposals for implementing blended educational supervision in Irbid Governorate from the perspective of educational supervisors?''}the researcher calculated the means and standard deviations for each item within the domain of developmental proposals. Table~\ref{tab8} presents the results.

{\small
\begin{longtable}[]{@{}
  >{\raggedright\arraybackslash}p{(\columnwidth - 10\tabcolsep) * \real{0.0513}} >{\raggedright\arraybackslash}p{(\columnwidth - 10\tabcolsep) * \real{0.6314}} >{\raggedright\arraybackslash}p{(\columnwidth - 10\tabcolsep) * \real{0.0683}} >{\raggedright\arraybackslash}p{(\columnwidth - 10\tabcolsep) * \real{0.0513}} >{\raggedright\arraybackslash}p{(\columnwidth - 10\tabcolsep) * \real{0.0512}} >{\raggedright\arraybackslash}p{(\columnwidth - 10\tabcolsep) * \real{0.1466}}@{}}
\caption{Means and Standard Deviations of Items Related to Developmental Proposals for Implementing Blended Educational Supervision (N = 180)}
\label{tab8}\\
\toprule\noalign{}
\endhead
\bottomrule\noalign{}
\endlastfoot
\textbf{No.} & \textbf{Item} & \textbf{Mean} & \textbf{SD} & \textbf{Rank} & \textbf{Level} \\
1 & Establishing comprehensive electronic infrastructure in education directorates and schools & 4.07 & 1.31 & 1 & High \\
7 & Appointing specialists in blended supervision within the educational supervision department & 4.06 & 0.69 & 2 & High \\
6 & Increasing financial resources to support supervisory methods through information networks & 4.05 & 0.31 & 3 & High \\
4 & Raising awareness among supervisors about the importance of using modern technologies in supervision & 4.03 & 2.64 & 4 & High \\
2 & Providing personal computers for educational supervisors & 4.02 & 0.85 & 5 & High \\
11 & Reducing administrative burdens on educational supervisors & 4.00 & 1.09 & 6 & High \\
10 & Conducting computer training courses for teachers & 3.98 & 0.97 & 7 & High \\
15 & Developing teachers' skills in using modern technologies & 3.96 & 0.95 & 8 & High \\
8 & Increasing the attention of the Educational Supervision Directorate toward blended supervision & 3.94 & 1.27 & 9 & High \\
13 & Enhancing teachers' motivation to use modern technologies & 3.93 & 1.13 & 10 & High \\
9 & Conducting computer training courses for educational supervisors & 3.90 & 1.10 & 11 & High \\
17 & Providing learning resources (books, guides, training manuals) specialized in blended supervision & 3.89 & 0.98 & 12 & High \\
12 & Enhancing supervisors' motivation to use modern technologies & 3.85 & 1.25 & 13 & High \\
14 & Developing supervisors' technical skills in dealing with new technologies & 3.80 & 1.01 & 14 & High \\
18 & Increasing financial and technical support for information network use by the Ministry of Education & 3.76 & 1.09 & 15 & High \\
16 & Providing the necessary software for remote supervision & 3.75 & 0.89 & 16 & High \\
3 & Providing personal computers for teachers & 3.74 & 0.90 & 17 & High \\
5 & Providing internet services in public schools & 3.63 & 1.13 & 18 & Moderate \\
\multicolumn{2}{@{}l}{\textbf{Overall Domain}} & \textbf{3.90} & \textbf{0.51} & -- & \textbf{High} \\
\end{longtable}
}

As shown in Table~\ref{tab8}, the mean scores for the developmental proposals ranged \textbf{between} 3.63 and 4.07, indicating that the overall level of agreement among respondents was high (M = 3.90,SD = 0.51)\textbf{.} The proposal with the greatest mean score was ``introducing a full electronics system in education directorates and schools'' (M = 4.07, SD = 1.31), whereas the least popular one was ``providing internet services in public schools'' (M = 3.63, SD = 1.13).

\section{Discussion}

The research investigated the extent to which conditions for the application of the blended educational supervision in Irbid region are being met, examined whether supervisors differ in their attitudes regarding the blended educational supervision as based on their demographic characteristics, and suggested possible amendments. The findings show that the EAAs indicated their strong readiness for blended educational supervision in the technical, material and human aspects and supported the practical recommendations. The implications of the findings indicate that educational supervision is moving from the initial stage of digital transformation to applying more comprehensive model of technological development in supervisory practice. However, the educational supervisors are still acknowledging the necessity of institutional support, continuous training, and considerable efforts and investments in technological infrastructure.

\subsection{Discussion of the First Research Question}

The first question was aimed at revealing why blended educational supervision has been made implementable in Irbid. The outcome showed a high degree of availability of all three dimensions, namely technical-electronic, material, and human requirements. The technical-electronic requirements received the most highly rated responses while material requirements came next in rank. The human requirements got the lowest means of all three requirements but still remained at a high level.

The results imply that educational supervisors perceive the main competencies necessary to implement blended supervision to be available. The high levels of technical and material requirements indicate that heavy investments have been made into Jordan education in the realm of digitization lately. It was stated that the recent growth of online learning practices and communication has made schools and educational directorates more skilled at using electronic systems, virtual communication, and technology in the sphere of administration, which can lead to better provision of the supervisory process.

In comparison to previous research conducted in Jordan and neighboring Arab countries, the present study reveals higher levels of readiness. Previous studies such as Hamdan \cite{ref21}, Hazaimeh \cite{ref18}, Salman \cite{ref22}, and Rajabi and Zamil \cite{ref23}generally identified moderate readiness levels, which accounts for limited technological capabilities, lack of professional development, and inconsistency of institutional support.

Higher levels of readiness in the present study can signify the development in the educational sector of Jordan due to the measures undertaken in the country aimed at promoting the digital transformation of education. This leads us to the conclusion that readiness for organizing blended supervising is not a static feature but something that varies depending on the advancement of technical environment and institutional competences.

Despite getting a high score on the scale, human needs are still the lowest in the structure. Based on the given facts we can conclude that technology readiness alone is insufficient for success. Supervisory staff still stresses the importance of development of professional competencies needed for digital supervising, joint leadership and proper application of technology in supervising. According to UNESCO \cite{ref1}, sustainable digital change in education requires not only technical capabilities but also constant professional education.

This research enriches the body of literature by demonstrating that the notion of blended supervision readiness should not be perceived merely as the availability of technology, but rather as a multidimensional process. The previous research has largely concentrated on the issues of technical or administrative readiness, whereas the current investigation shows that the supervisors consider human, material, and technological needs as interrelated components of efficient supervision in Jordanian schools.

\subsection{Discussion of the Second Research Question}

\noindent The second research question focused on whether the supervisors' perceptions were different based on the article's gender, and previous experiences as well as levels of education. The results indicated that the analyzed demographic characteristics did not significantly affect the supervisors' perceptions of the need for blended supervision. In fact, the lack of demographic differences suggests that the supervisors' awareness of blended supervision is more or less homogenous among supervisors, who are uniformly influenced by the system of performance standards in their work that is regulated by the Ministry of Education in Jordan.

\noindent This conclusion is corroborated by the findings of previous and similar research conducted by Hazaimeh \cite{ref18}, and Salman \cite{ref22}, because these authors reported that the demographic characteristics have negligible influence on the supervisors' perceptions of blended supervision in the Arab countries. It appears that the results of Hamdan \cite{ref21} show the same tendency because the researcher stated that gender as well as years of experience do not affect the readiness of the supervisors for supervision with the help of technological devices and methods. This conclusion is corroborated by the findings of previous and similar research conducted by Hazaimeh \cite{ref18}, and Salman \cite{ref22}, because these authors reported that the demographic characteristics have negligible influence on the supervisors' perceptions of blended supervision in the Arab countries. It appears that the results of Hamdan \cite{ref21}, show the same tendency because the researcher stated that gender as well as years of experience do not affect the readiness of the supervisors for supervision with the help of technological devices and methods.

\noindent In practical terms, the lack of statistically significant differences has important implications. In practical terms, this means that future professional development programs will not be required to target specific demographic groups. It is possible to develop the national training curriculum aimed at enhancing blended supervision skills within the entire supervisory workforce. This allows for more consistency in supervision practice and promotes the use of blended supervision at various educational directorates.

\subsection{Discussion of the Third Research Question}

\noindent The third question in this study looked into the proposals for development essential for the efficient use of blended educational supervision practices. The supervisors agreed greatly with all the measures suggested, particularly with those relating to enhancing the digital capacities, hiring the specialists in blended supervision, providing sufficient funding, and enhancing the professional training of supervisors and teachers. Besides, it shows that supervisors regard blended supervision as more than just the implementation of digital technologies. In their view, blended supervision is a process that involves a combination of improvements in technology, human resources, organizational support, and professional development. The findings confirm the fact that successful educational innovations depend on the combination of technology and improvements in human resources and institutional policy.

\noindent As for the priorities recognized in the research, they are more or less in line with those identified in the preceding studies. Hamdan \cite{ref21}, Rajabi and Zamil \cite{ref23} and Mohamed, Hassan, and Salim \cite{ref15} focus, in particular, on the significance of digital technologies, professional development, and institutional support. However, in this research, the contribution is made to the previous studies by showing that the supervisors themselves give priority to the mentioned directions.

\noindent One important thing to note is that internet services showed the lowest rating among the presented developmental options, even though the overall level of rating was moderate. It indicates that supervisory specialists still consider the issue of internet connectivity a domain where there is a demand for further development despite the positive changes made in the current environment of digital transformation. Reliable internet connectivity is crucial for synchronous communication, watching virtual classes, giving feedback online, and cooperative work, which explains why the issue of internet connectivity development has to be dealt with further for blended supervision implementation in schools and educational directors' offices.

\noindent In general, the developmental recommendations suggested in the study can be useful for practitioners in educational policymaking. Rather than just paying attention to technology-related purchases, supervisors suggested the use of a comprehensive solution.

\section{Contribution of the Study}

Research on blended educational supervision has been increased by including empirical evidence from Jordan, where literature on this topic seems scarce. Different from a number of previous studies that mostly dealt with levels of technological readiness, the present paper attempted to examine the combination of technological, material and human requirements and represent formed priorities by educational supervisors. Results of the study indicate that blended educational supervision has to be put into practice through the combination of technology, skills and support available rather than through any of these aspects separately. The findings provide evidence that can contribute to future policy-making and help the leaders in education improve blended supervision approaches.

\subsection{Theoretical and Practical Contributions}

This study enriches the existing literature with recent empirical data from the Jordanian educational context, which has been sparse when it comes to studies focusing on this specific area of study. Instead of focusing solely on a single requirement such as human, technical, and material, this study explores the overall picture by discussing all three requirements simultaneously. Moreover, we explore some proposed developmental ideas suggested by educational supervisors regarding blended supervision. Additionally, it provides an extensive knowledge of how the different human, technical, and material needs might be satisfied in practice for the successful execution of blended supervision at the school level.

Moreover, the current study makes another contribution to the regional literature through providing empirical evidence coming from Irbid governorate utilizing both descriptive and inferential statistics. Whereas most studies conducted in Jordan and in other Arab countries focused on the description of the existence of blended or electronic supervision needs, the current study examines whether the perception of educational supervisors differs based on their gender, professional experience, and educational qualifications. It turns out that there are no statistically significant differences between those variables suggesting that perceptions about blended supervision become relatively uniform among educational supervisors, probably due to common institutional policies and professional practices.

Supervisors highlighted several priorities for development including enhancing the country's digital infrastructure, increasing professional development opportunities, improving financial support, and employing specialists in blended supervision. In conclusion, this study provides evidence-based suggestions that support existing initiatives to promote digital innovations across various sectors in Jordan while providing a basis for future organizational changes that should be implemented by educational leaders.

The broader implications of our results also confirm a central insight from prior work indicating that success in blending technology into an educational supervisory system relies upon mutual interactions among various factors related to technological resources, professional competencies, and institutional support , Furthermore, these results complement earlier regional studies by emphasizing the need for coordinated investments in building organizational capacities alongside human development to drive long-term sustainability of educational supervisory reforms.

Finally, more generally, our findings suggest continued validity for increasingly widespread recognition that effective educational supervision entails creating a dynamic process where educators collaboratively engage in continuous professional learning experiences grounded in reflection and educational enhancement via integrated use of digital technologies. By providing such context specific evidences, this current study lays important groundwork for future scholarship and/or policies regarding educational supervisory developments particularly within countries sharing similar education systems such as that in Jordan itself.

\section{Conclusion}

This study explored the availability of the requirements needed for the implementation of BES in Irbid governorate from the perspectives of the educational supervisors (N=61) and outlined some possible ways forward which could help achieve success. It was found out that there was high availability in terms of Technical-electronic, material and human resources. In other words; educational supervisors perceived that their organizations are equipped with all the necessary conditions required for successfully implementing BES. Also, there was a general acceptance among supervisors towards various types of developmental interventions such as enhancing digital infrastructures, giving more chances to enhance professional growth and increase the number of institutional supports through involving people with special skills regarding BES.

Finally, no statistically significant differences were found when it comes to supervisor's perception based on gender, professional experiences and educational qualifications respectively. Hence, these findings imply that ensuring successful implementation of blended educational supervision relies mainly on integrating technology tools, offering continuous education and building up an environment with enough internal and external support. Therefore, by providing empirical evidence on blended educational supervision from the case of Irbid governorate, this work adds value to the small number of studies conducted on blended educational supervision in Jordan. Moreover, the results offer solid proof which can be used by educational leaders and policy makers while developing plans about future supervisory reforms aiming at digital transformations within the educational sector.

\subsection{Recommendations}

Based on the findings of this study, the following recommendations are proposed:

\begin{itemize}
\item
  \begin{quote}
  Continue investing in digital infrastructure, including hardware, software, communication networks, and online supervisory platforms to ensure the effective implementation of blended educational supervision.
  \end{quote}
\item
  \begin{quote}
  Expand continuous professional development programs that strengthen supervisors' and teachers' digital, pedagogical, and supervisory competencies.
  \end{quote}
\item
  \begin{quote}
  Develop institutional policies and administrative procedures that support the integration of blended supervision into routine supervisory practice.
  \end{quote}
\item
  \begin{quote}
  Appoint specialists in blended educational supervision to provide technical guidance and professional support for supervisors and schools.
  \end{quote}
\item
  \begin{quote}
  Encourage the exchange of successful supervisory practices among educational directorates to promote continuous improvement and innovation.
  \end{quote}
\item
  \begin{quote}
  Extend future research to other governorates in Jordan and examine the relationship between blended educational supervision and variables such as instructional quality, professional development, organizational innovation, and job satisfaction.
  \end{quote}
\end{itemize}

The effective application of blended learning supervision necessitates a synchronized national strategy that integrates technological advancement, ongoing professional development, and conducive institutional policies. Improving these aspects will enable the development of a long-lasting supervision system that can ensure quality education and promotes digital transformation in Jordanian schools.

\textbf{Ethical considerations}

The research was undertaken in line with the ethical guidelines established for research involving human beings. It was conducted with the approval of Jadara University, Faculty of Education Ethics Committee. The approved protocol number for this research is (to be inserted when officially assigned). Participation in the research was voluntary. The educational supervisors who participated were intimated about the aims and objectives of the research before data collection. Written informed consent was received from all the participants, making it easy for them to drop out at any stage of research without any problems.

\textbf{Conflict of Interests:} The author declares that no conflict of interest concerning the submission of this research exists.

\textbf{Funding:} This research did not receive any specific grant from funding agencies.

\textbf{Use of Artificial Intelligence}

Artificial Intelligence was used marginally in the process of editing and enhancing the language used in the research. The interpretation of data, and the entire research design and intellectual content were the responsibility of the author. The author guarantees the originality and authenticity of the work.

\textbf{Acknowledgments:} We would like to express our gratitude to \textbf{Jadara} University for recognizing the researchers' efforts in completing this research\textbf{.}

\subsection{}

\begin{thebibliography}{99}

\bibitem{ref1} UNESCO. (2022). \emph{ICT competency framework for teachers}. UNESCO.

\bibitem{ref2} Simonson, M., \& Seepersaud, D. (2018). \emph{Distance education: Definition, history, and future directions}. Information Age Publishing.

\bibitem{ref3} Vaiz, F., Rahman, L., \& Nordin, M. (2021). \emph{Integrating technology in educational supervision: A comparative study}. International Review of Education, 67(4), 493--510.

\bibitem{ref4} Mohamed, A., Hassan, F., \& Salim, K. (2023). \emph{Blended educational supervision: Status, enablers, and constraints in Azhar secondary institutes}. Journal of Educational Practice, 14(1), 55--74.

\bibitem{ref5} Hussein, A., \& Al-Qathami, M. (2019). The role of technology in modern supervisory models. \emph{Arab Journal of Educational Research, 9}(4), 45--61.

\bibitem{ref6} Bani Younes, Z. B. (2025). Investigating how AI algorithms can be used to automate English language proficiency assessments. World Journal of English Language, 15(5), 285. https://doi.org/ 10.5430/wjel. v15n5p285

\bibitem{ref7} Salihuddin, S. (2020). \emph{The role of digital readiness in supporting blended educational supervision}. International Journal of Educational Technology and Learning, 8(3), 45--57.

\bibitem{ref8} Salihuddin, N. (2020). \emph{Efficiency and innovation in blended supervisory systems}. Asia-Pacific Journal of Educational Technology, 12(1), 33--49.

\bibitem{ref9} Arnauld, E. M. (2016). \emph{Virtual supervising: Perceptions of interaction and feedback during student teaching} (Doctoral dissertation, University of New Mexico). ProQuest Dissertations \& Theses Global (UMI No. 10100453).

\bibitem{ref10} Al-Amri, F. (2023). \emph{The reality of practicing blended educational supervision in secondary schools in Al-Baha City}. Journal of Bisha University for Educational Sciences, 6(2), 744--778.

\bibitem{ref11} Al-Mubayyid, A. (2014). \emph{Educational supervision and teacher development in contemporary contexts}. Amman: Dar Al-Fikr.

\bibitem{ref12} UNESCO. (2023). \emph{Digital transformation of education: Policy framework and implementation guide}. UNESCO.

\bibitem{ref13} Bani Younes, Z., Talafha, D. K., \& Gul, K. (2025). Artificial intelligence as a writing coach: Impacts on BS English students. Jordan Journal of Modern Languages \& Literatures, 17(4), 891--918. https://doi.org/10.47012/jjmll.17.4.3

\bibitem{ref14} Zhong, L., Vasconcelos, L., Cheng, M., \& Jackson, S. (2019). Leadership development through organizational efforts: Reflections by interns at the 2018 Association for Educational Communications and Technology International Convention. \emph{TechTrends, 63}(6), 650--651.

\bibitem{ref15} Mohamed, A., Hassan, N., \& Salim, M. (2023). \emph{Requirements for implementing blended supervision in Al-Azhar secondary institutes}. Helwan University Journal of Educational and Social Studies, 29(2), 140--175.

\bibitem{ref16} Alharthi, S., \& Alghamdi, M. (2021). \emph{Collaborative supervision and teacher development in blended contexts}. Journal of Educational Supervision, 34(2), 98--115.

\bibitem{ref17} Bani Younes, Z. B. (2025). Computer-assisted language learning: Insights from Jordanian private university English as a foreign language undergraduate and instructors. Forum for Linguistic Studies, 7(1), 355--368. https://doi.org/10.30564/fls.v7i1.7807

\bibitem{ref18} Al-Hazaymeh, A. (2020). \emph{The availability of e-supervision requirements in the northern governorates of Jordan from the perspective of educational supervisors}. Jordanian Educational Review, 45(3), 251--268.

\bibitem{ref19} Hamdan, A., \& Al-Ajeez, M. (2015). \emph{Technological integration in educational supervision: A framework for reform}. Journal of Instructional Leadership, 25(4), 203--220.

\bibitem{ref20} Al-Hanai, M. (2017). \emph{Challenges of traditional supervision in Arab educational systems}. Arab Journal of Educational Studies, 12(1), 33--47.

\bibitem{ref21} Al-Hamdan, M. (2015). \emph{The degree of availability of e-supervision requirements in public schools in Gaza governorates and ways to develop them} (Unpublished master's thesis). Faculty of Education, Islamic University of Gaza.

\bibitem{ref22} Salman, M. (2021). \emph{The reality of blended supervision practices in public schools and ways of developing them in light of modern trends from the perspectives of educational supervisors in the central directorates of the northern governorates} (Unpublished master's thesis). Faculty of Education, Al-Quds Open University, Palestine.

\bibitem{ref23} Rajabi, Y., \& Zamil, M. (2022). \emph{A proposed model for blended educational supervision in Palestinian public schools from the perspectives of school principals and educational supervisors}. Al-Quds Open University Journal for Educational and Psychological Research and Studies, 13(39), 181--199.

\bibitem{ref24} Al-Ghamlas, A. (2008). \emph{Foundations of educational supervision in contemporary contexts}. Riyadh: Dar Al-Mareekh.

\bibitem{ref25} Schwartz-Bechet, B. (2014). \emph{Virtual supervision for teachers and supervisors: Effectiveness and feasibility of ICT use in supervision}. International Journal of E-Learning and Teaching, 10(4), 221--238.

\bibitem{ref26} Arnauld, E. M. (2016). \emph{Virtual supervising: Perceptions of interaction and feedback during student teaching} (Doctoral dissertation, University of New Mexico). ProQuest Dissertations \& Theses Global (UMI No. 10100453).

\bibitem{ref27} Hamdan, A. (2015). \emph{Availability of e-supervision requirements in Gaza government schools and ways to develop them}. Journal of Instructional Leadership, 25(4), 203--220.

\bibitem{ref28} Al-Juhani, M., \& Al-Juhani, R. (2023). \emph{The role of supervisors across the three stages of blended supervision and requirements for its implementation in general education in Madinah}. Journal of Educational Research and Practice, 15(2), 145--170.

\bibitem{ref29} Banga, H., \& Al-Qathami, M. (2019). \emph{A proposed program for implementing a blended supervision model from the perspective of teachers and educational supervisors in light of modern trends}. Arab Journal of Qualitative Education, 10(3), 83--112.

\bibitem{ref30} Mohamed, A., Hassan, F., \& Salim, K. (2023). \emph{Blended educational supervision: Status, enablers, and constraints in Azhar secondary institutes}. Journal of Educational Practice, 14(1), 55--74.

\bibitem{ref31} Creswell, J. W., \& Creswell, J. D. (2023). \emph{Research design: Qualitative, quantitative, and mixed methods approaches} (6th ed.). SAGE.

\bibitem{ref32} Field, A. (2021). \emph{Discovering statistics using IBM SPSS Statistics} (6th ed.). SAGE.

\end{thebibliography}

\end{document}
